
\section{The State, Taxes and Tariffs}

\begin{frame}{Flow Model of the Economy}
    \showmytikz[1,2,3,4]{../shared/tikz/tikz-flowmodel-own}
    \mycomment{
    The Economy is modelled as consisting of consumers, firms (companies, businesses) and the state.

    However, people have different roles simultaneously: as consumers, but also act as firms. 

    The government, however, may be democratic or not, and will pursue an agenda for itself, for consumers, and for firms; interests that may compete with one another.
    }
\end{frame}


\begin{frame}{Types of Policy}
    \mytext{
    \begin{itemize}
        \item Fiscal Policy
        \begin{itemize}
            \item Government spending
            \item Taxation, tariffs
            \item Borrowing
        \end{itemize}
        \item Regulatory Policy
        \begin{itemize}
            \item restrictions
            \item obligations
            \item standards
        \end{itemize}
        \item Monetary Policy
        \begin{itemize}
            \item interest rates
            \item open market policy
        \end{itemize}
    \end{itemize}
    }
\end{frame}
 

% \begin{frame}{Policy Actors: Private Sector}
%     \mytext{
%         \begin{itemize}
%             \item consumers: workers \& employees
%             \item producers: firms \& corporations
%             \item traders: merchands, transporters
%             \item banks, financial institutions
%             \item port authorities, associations
%         \end{itemize}
%     }   
% \end{frame}

\begin{frame}{Levels and Institutions of the State}
    \mytext{
        \begin{itemize}
            \item Government
            \begin{itemize}
                \item Federal level
                \item Länder level
                \item Communal level
                \item Social Security
                \item Governmental Corporations
                \item Supranational level (EU)
            \end{itemize}
            \item Central Bank
        \end{itemize}
    }
    \mycomment{
    A typical modern state has three branches: the executive, judicative and legislative branch. On top, there are various levels of government, particularly in federal states and in the European Union.
    
    In counties with Statutory Law (as opposed to Common Law), private sector entities act under Private Law (\textit{Privatrecht}), while public sector entities fall under Public law (\textit{Öffentliches Recht}).

    Similarly, consumers may legally be treated differently than companies (e.g. \textit{BGB, HGB}) and corporations.
    }
\end{frame}



\begin{frame}{Shell Model}
    \showmytikz{../shared/tikz/tikz-shellmodel}
    \mycomment{
        The `shell model' was developed in the framework of the European System of Accounts (ESA; German: \textit{ESVG}), so that budgets can be compared. 

        Governments may move parts of the budget into extra budget for cosmetic or practical reasons. \cite{destatisbegriffe2019}

        Governments can levy taxes, and - in general - hold a monopoly on the use of force.
    }
\end{frame}



\begin{frame}{Government Objectives}
    \mytext{ 
    Governments may care for...
    \begin{itemize}
        \item (re-)elections, welfare, influence, popularity, income
        \item income distribution, equality, infrastructure, education, growth, tax base, exports, state monopolies, foreign relations, sustainability, government revenue and tax income, money supply \& demand, debt level, public safety, defense ... 
        % \item fiscal policy, regulatory policy; selling sovereign debt
    \end{itemize}
    }
    \mycomment{
    There are multiple level of government, and quite often it is not as clar-cut what is governmental and what is private.
    }
\end{frame}


\begin{frame}{Government Expenditure Ratio}
    \mytext{
       The state is a powerful actor. Just how large the role of the state in a country is can be measured by the state quota, that is the government spending ratio. 

        \begin{equation}
            SQ=\frac{G}{Y}\%
        \end{equation}

        \mytable{
            \begin{tabular}{lr}
                \toprule
                Ukraine & 66\%\\
                    France & 58\% \\
                 Germany & 49\%  \\
                 Japan & 44\%\\
                 USA &  36\% \\
                 Saudi Arabia & 28\%\\
                 Pakistan &  20\% \\
                 \bottomrule
            \end{tabular}
        }
    }
        
    \mycomment{
    The government spending ratio varies widely, with low income countries typically lower in the table. The highest ratio are offshore havens and other outliers. War efforts obviously increase this ratio.
    }
\end{frame}


\begin{frame}{The Central Bank}
    \mytext{
    An important part of the State is the central bank:
    \begin{itemize}
        \item cares for / objectives (?):
        \begin{itemize}
            \item profits
            \item price stability
            \item currency / exchange rate stability
            \item asset stability 
        \end{itemize}
        \item tools:
        \begin{itemize}
            \item money / currency supply
            \item via interest rates
            \item bank regulation, minimum reserve
            \item credit and debit market regulations
            \item exchange rate policy
            \item open market policy: buying private and government debt
        \end{itemize}
    \end{itemize}
    }
    \mycomment{
    The central bank is a part of government, albeit a very special one, which stems from the privilege of the State to issue currency. 

    The central bank is the main influencer of capital markets. It may also decide to buy government bonds, thus financing government debt. Buying bonds on the open market can provide the economy with funds (quantitative easing). In case of crises, the central bank acts as lender of last resort.

    Most central banks today act fully independent of the government. Their own interest and that of the economy should serve the economy best.
    }
\end{frame}



% \begin{frame}{Pareto Optimum, Pareto Criterium}
%     \mytext{
%     A social situation is a Pareto optimum, if nobody can be made better off without anyone else being made worse off.
%     \begin{itemize}
%         \item exchange leads to Pareto optimum
%         \item perfect markets lead to Pareto optimum
%     \end{itemize}
%     Pareto criterium:
%     \begin{itemize}
%         \item if a situation is not Pareto-optimal, but a Pareto optimum can be achieved, act so that the optimum gets achieved
%     \end{itemize}
%     }
%     \mycomment{
%     The Pareto criterium acts as a minimalist criterium. All people will always prefer situations that makes everyone better off.

%     The exchange can also be non-materialistic; I might just be better off a little poorer in monetary terms and not having to watch my neighbour suffer. In this case, a voluntary donation moves us to Pareto optimum.
%     }
% \end{frame}

% \begin{frame}{Kaldor-Hicks criterium}
%     \mytext{
%     Improving some $A$ at the cost of $B$ (violating the Pareto criterium) can be justified when some of the gains of $A$ can be redistributed (later) to $B$.

%     \begin{itemize}
%         \item redistribution by tax and transfer
%         \item redistribution by growth and future gains (of $B$)
%     \end{itemize}
%     }
%     \mycomment{
%     The Kaldor-Hicks criterium is similar to the Pareto criterium, but it requires a mechanism for the redistribution. 
%     }
% \end{frame}

% \begin{frame}{Pareto Optimum}
%     \showmytikz{../shared/tikz/tikz-pareto.tex}
%     \mycomment{
%     A \textbf{Pareto optimum} is the situation where no improvement to one party can be made without leaving the other party worse off.
    
%     In this graph the utility of of a good allocated between two parties are plotted as $X,Y$ coordinates. 

%     All points on the graph are a Pareto optimum, as there is no other point for which both $A$ and $B$ would be better off. This line is called \textbf{Pareto frontier.}

%     The \textbf{Kaldor-Hicks} criterium states that even a non-Pareto-optimal situation may be desirable if and only if the disadvantage for the losing party could be compensated in another transaction. For example, a large win for $A$ and a small loss for $B$ could be compensated by a transfer, e.g. by taxing $A$ and benefitting $B$.
%     }
% \end{frame}


% \begin{frame}{Political Economy - Descriptive}
% \mytext{
% Actual motivation of the state?
%     \begin{itemize}
%         \item enrichment / ambition of elite?
%         \item benevolent / ideology?
%         \item equalisation or growth?
%         \item income of state owned enterprises?
%         \item inland revenue?
%         \item rational at all?
%         \item tradition? re-election?
%     \end{itemize}
% }
% \end{frame}


\begin{frame}{External Effects}
    \showmytikz[1]{../shared/tikz/tikz-externalities-pigou}
    \mycomment{
    Whatever we do, we typically incur external effects. We are not paying for all consequences of our actions; we are using common goods without paying; and even if we are forced to pay, some damages are not redeemable by payment.
    }
\end{frame}

\begin{frame}{External Effects}
    \showmytikz[1,2]{../shared/tikz/tikz-externalities-pigou}
    \mycomment{
    Quite often, the costs of the producer do not cover all costs. Those costs not covered are called externalities, or social costs.

    A classic example is the pollution caused by production of the good, or in fact any other type of damage or costs.

    These are hard to grasp, and are often only realised much in the past or by future generations. We abstract from those problems and assume $MC_{social}$ is known.
    }
\end{frame}

\begin{frame}{External Effects}
    \showmytikz[1,2,3]{../shared/tikz/tikz-externalities-pigou}
    \mycomment{
        Social costs are above the costs covered by the producers and the price the consumer pay them. Some of them are offset by the consumer and producer surplus: yes, we destroy the planet, but hey, we consume!

        Bear in mind that the consumer realising the surplus may not be the same consumers suffering the damage!

        However, there is an excess damage that is neither translated into consumer nor producer surplus. This is sometimes called dead weight loss, but it is actually excess costs, or non-covered damage. (Dead weight loss is welfare we could have but do not have due to market imperfections. Here, we actually have a loss, so it is not "dead weight").
    }
\end{frame}

\begin{frame}{External Effects}
    \showmytikz[1,2,3,4]{../shared/tikz/tikz-externalities-pigou}
    \mycomment{
        The idea could be to internalise the social costs, lifting the cost curve (supply curve) up to meet the social cost curve. Then the price realised would be realistic. This could be done by a tax.

        This tax can achieve two things at once:
        \begin{itemize}
            \item shifting the price so that private costs equal social costs
            \item raising a government revenue that can be used to compensate those who suffer without benefiting.
        \end{itemize}
    }
\end{frame}


% \begin{frame}{Median Voter Theorem}
%     \showmytikz{../shared/tikz/tikz-medianvoter}
% \end{frame}

\begin{frame}{Median Voter Theorem}
    \showmytikz[1,2,3]{../shared/tikz/tikz-medianvoter-two}
    \mycomment{Any majority must include the `guy in the middle', the median voter.}
\end{frame}


\begin{frame}{Taxing a Market}
    \showmytikz[1,2,3,4]{../shared/tikz/tikz-tax.tex}
    \mycomment{
        The graph shows the standard microeconomic analysis of taxes applied to a market. The tax shifts the supply curve up, reducing both consumer and producer surplus. In addition, there is a dead weight loss. 
        
        Who bears the larger share of the welfare loss---and thus the weight of the tax---depends on the elasticities of demand and supply, rather than political statements. In particular, it does not matter if producers or suppliers pay the tax.
    }
\end{frame}

\begin{frame}{Laffer Curve}
    \showmytikz[1,2]{../shared/tikz/tikz-laffer}
    \mycomment{
    The Laffer curve illustrates that 
    \begin{itemize}
        \item at a tax rate of 100\%, nobody would work anymore, and government revenue would be zero;
        \item at a tax rate of 0\% nobody would have to pay any taxes, so government revenue would also be zero.
    \end{itemize}
    The exact curve is highly disputed. 
    }
\end{frame}


\begin{frame}{The Open Economy}
    \showmytikz[1,2]{../shared/tikz/tikz-flowmodel-open}
    \mycomment{
    We understand the economy of a country to consist chiefly of labour markets, the markets for goods and services, and capital markets.

    The state influences in various ways. It taxes its own citizens, regulates activities, provides currency, and purchases from them.
    }
    
\end{frame}

\begin{frame}{The Open Economy}
    \showmytikz[1,2,3]{../shared/tikz/tikz-flowmodel-open}
    \mycomment{
    But the economy is not alone. There are foreign countries, and there is indeed a world market, where there are also markets for labour, goods and services, and capital. 
    
    Because goods can easily be re-routed if there is arbitrage, in effect there is a highly efficient world market. 
    }
    
\end{frame}

\begin{frame}{The Open Economy}
    \showmytikz[2,4]{../shared/tikz/tikz-flowmodel-open}
    \mycomment{
    Unless there are barriers, and abstracting from transportation costs, the markets for goods and services merge. 

    The markets for labour may or may not merge; but labour is clearly less mobile than goods and capital, due to administrative and language barriers, but also due to family bounds.

    There is, however, a connection here too: firms may decide to produce abroad, investing abroad (foreign direct investment, FDI), and (re-)import.

    Capital markets were not always free, but there are no more capital restrictions at present. Monetary flows are also not taxed. The crossing of goods (and services) across borders might be taxed using tariffs or quotas.
    }
\end{frame}

\begin{frame}{Governance of the Open Economy}
    \mytext{
        With trade, the tools of the state become less effective. 
        
        The state may want to influence the outcome. 

        \begin{itemize}
            \item controls at the border
            \begin{itemize}
                \item migration (labour)
                \item capital (?)
                \item goods and services
            \end{itemize}
            \item tariffs, quotas
            \item regulation, other barriers
            \item subsidies
            \item tax competition
            \item monetary policy
            \item contracts, diplomacy
        \end{itemize}
    }
    \mycomment{
        In an open economy, all government actions have an impact not only on the economy, but also abroad; or it may become less effective or ineffective in an open economy. This depends on the (relative) size of the economy (or state). 

        The state can direct policy directly to the cross-border flows, such as tariffs. But also policy that is directed at the "inside" such as taxes, regulations or subsidies, will effect cross-border flows.
    
        % Policies to influence market outcomes are typically limited in scope. Sometimes, they may not be successful at all, have detrimental effects, or side effects.
    }
\end{frame}

% \section{Protectionism I: Tariffs and Quotas}

\begin{frame}{Effect of Tariffs}
    \showmytikz[1,2]{../shared/tikz/tikz-tariff.tex}
    \mycomment{
    In a country without trade (and tax), local demand is satisfied only by local supply.
    
    In an open economy, importers will offer the good at world market price. Thus only local producers that are cheaper than world market producers from abroad will still serve the local market. 
    }
\end{frame}

\begin{frame}{Effect of Tariffs}
    \showmytikz[1,2,3]{../shared/tikz/tikz-tariff.tex}
    \mycomment{
        The import has resulted in a new supply curve leading up to the world price, when it becomes flat, that is: effectively infinitely elastic.
        
        (Local) producer surplus is lower, while consumer surplus is higher due to larger quantity and lower price.
    }
\end{frame}

\begin{frame}{Effect of Tariffs}
    \showmytikz[1,2,3,4]{../shared/tikz/tikz-tariff.tex}
    \mycomment{
        Now the government can raise a tariff. Every unit imported costs now more, effectively rising the local price from world market price to world market price plus tariff. 
    }
\end{frame}

\begin{frame}{Effect of Tariffs}
    \showmytikz[1,2,3,4,5]{../shared/tikz/tikz-tariff.tex}
    \mycomment{
        This benefits local producers, while it rises the local price and decreases consumer surplus.
    }
\end{frame}

\begin{frame}{Effect of Tariffs}
    \showmytikz[1,2,4,5,6]{../shared/tikz/tikz-tariff.tex}
    \mycomment{
        The tariff will raise some government revenue (GR), which the government can redistribute. 
        
        However, collecting the tariff will also come at a cost.

        Any tariff will result in dead weight loss (DWL) that cannot be redistributed, as it is simply lost. This burden is carried by the consumer, as it is consumer surplus that is lost.

        Notice that there is a certain tariff that maximises government revenue.
    }
\end{frame}

\begin{frame}{Effect of Tariffs}
    \showmytikz[1,2,4,5,6,7]{../shared/tikz/tikz-tariff.tex}
    \mycomment{
        The tariff has reduced the import quantity, while the world market price has not changed; but consumers pay a higher price due to the tariff. 
        
        So in total consumers may spend just as much as before on the imported goods, albeit for a smaller quantity (compare the red and the green rectangle). Foreign producers sell less; but the burden of the tariff is on the consumer (GR).
    }
\end{frame}

\clearpage

\begin{frame}{Effect of Quotas}
    \showmytikz[1,2]{../shared/tikz/tikz-quota.tex}
    \mycomment{
        An import quota restricts the amount allowed to enter the country.

        Note that export restrictions are exactly similar, only that the "quota" is then imposed by the exporting country.
    }
\end{frame}

\begin{frame}{Effect of Quotas}
    \showmytikz[1,2,3]{../shared/tikz/tikz-quota.tex}
    \mycomment{
        This effectively reduces the volume traded. The welfare effect can be the same as a tariff, only that the quota does not raise government revenue.
    }
\end{frame}

\begin{frame}{Effect of Quotas}
    \showmytikz[1,2,3,4]{../shared/tikz/tikz-quota.tex}
    \mycomment{
        However, how quotas are distributed remains in question.

        The government could use a lottery, a first-come-first-served system, or auction the quotas. If the quota do not cost the importers, they are effectively a present to them (importer surplus, $IS$).

        We see: a quota can have the same economic effect as a tariff. In the case of tariff, the price is manipulated; in the case of quota, the quantity. Since price and quantity depend on one another, the result can be the same.
    }
\end{frame}


\begin{frame}{Trump "Reciprocal" Tariffs}
\mytext{
    On April 2nd 2025, US President Trump presented a table with high tariffs imposed on different countries. The underlying formula was \cite{trumpstariffscalculation}:

    \begin{equation}
        \Delta\tau_i=\frac{x_i-m_i}{\varepsilon\cdot\varphi\cdot m_i} 
    \end{equation}

    \begin{description}
        \item [$x_i-m_i$] trade imbalance with the US
        \item [$\varepsilon$] import elasticity, set to $\varepsilon=4$ (\textit{while maintaining} $\varepsilon<0$)
       \item [$\varphi$] passthrough from tariffs to import prices, set to $\varphi=0.25$
    \end{description}

    Which becomes:

    \begin{equation}
        \Delta\tau_i= \frac{m_i-x_i}{2\cdot m_i}
    \end{equation}

    \mycomment{
    This is like saying:
    \begin{quote}
        "If what I sell is worth half as much as what I buy,\\ let me pay 25\% more on what I buy---particularly when I don't buy much from there in total."
    \end{quote}

    Imagine my daughter to spend more than she earns. Will it better the situation if I double the price for everything she buys?

    This does indeed depend on the elasticities of exports (what she sells) and imports (what she buys). The US is chiefly importing goods, and exporting liabilities.
   
    }
}

\end{frame}

